\section*{Conclusão}

Este trabalho apresentou uma análise comparativa e padronizada dos algoritmos de agrupamento KMeans, Fuzzy C-Means e Agglomerativo aplicados aos conjuntos de dados Adult e Dry Bean. A padronização do pipeline experimental, com balanceamento das classes, repetição dos experimentos, uso de métricas quantitativas e visualizações, permitiu uma avaliação justa e didática do desempenho de cada método.

Os resultados evidenciaram que:
\begin{itemize}
    \item O \textbf{KMeans} apresentou desempenho consistente, com boa acurácia média e baixo desvio padrão, sendo eficiente para clusters bem separados e esféricos.
    \item O \textbf{Fuzzy C-Means} mostrou-se útil em cenários com sobreposição entre grupos, permitindo análise de pertinência, mas pode ser mais sensível a ruídos e à escolha do número de clusters.
    \item O \textbf{Agglomerativo} destacou-se pela flexibilidade em capturar estruturas hierárquicas, mas pode ser mais afetado por outliers e pela definição da métrica de ligação.
\end{itemize}

A tabela a seguir resume a comparação dos métodos quanto a desempenho médio e robustez (desvio padrão), apresentando lado a lado os resultados quantitativos para Adult e Dry Bean:

A seguir, apresentamos tabelas separadas para cada conjunto de dados, detalhando a acurácia média e o desvio padrão dos algoritmos avaliados.

\begin{table}[H]
\centering
\caption{Comparação dos algoritmos de agrupamento no conjunto Adult}
\begin{tabular}{lcc}
\toprule
\textbf{Algoritmo} & \textbf{Acurácia Média} & \textbf{Desvio Padrão} \\
\midrule
KMeans & 0,8221 & 0,0142 \\
Fuzzy C-Means & 0,8010 & 0,0175 \\
Agglomerativo & 0,6932 & 0,0148 \\
\bottomrule
\end{tabular}
\label{tab:comparacao_adult}
\end{table}

\begin{table}[H]
\centering
\caption{Comparação dos algoritmos de agrupamento no conjunto Dry Bean}
\begin{tabular}{lcc}
\toprule
\textbf{Algoritmo} & \textbf{Acurácia Média} & \textbf{Desvio Padrão} \\
\midrule
KMeans & 0,7215 & 0,0123 \\
Fuzzy C-Means & 0,6840 & 0,0157 \\
Agglomerativo & 0,6125 & 0,0112 \\
\bottomrule
\end{tabular}
\label{tab:comparacao_drybean}
\end{table}

De modo geral, a escolha do algoritmo mais adequado depende da natureza dos dados, do objetivo da análise e da necessidade de interpretabilidade ou robustez frente a ruídos e desbalanceamento. O estudo reforça a importância de pipelines padronizados e análise crítica para a aplicação prática de técnicas de agrupamento em ciência de dados.
